\documentclass[12pt,a4paper]{article}
\usepackage[margin=2cm]{geometry}
\usepackage{graphicx}
\usepackage{tabularx}
\usepackage{array}
\usepackage{enumitem}
\usepackage{hyperref}
\usepackage{float}
\usepackage{fontspec}
\defaultfontfeatures{Ligatures=TeX}
\usepackage{tikz}
\usetikzlibrary{calc}

\setlength{\parindent}{0pt}
\setlength{\parskip}{6pt}

% A simple checkbox
\newcommand{\checkbox}{\(\square\)}

% A line field
\newcommand{\linefield}[1]{\rule{#1}{0.4pt}}

% Task heading
\newcommand{\tasktitle}[2]{\vspace{6pt}\textbf{#1}\hfill /#2\par\vspace{4pt}}

\usepackage{amssymb}
\providecommand{\checkbox}{$\square$}

\begin{document}

% ===== HEADER FRAME =====
    \noindent
    \setlength{\fboxsep}{8pt}   % Innenabstand im Rahmen
    \setlength{\fboxrule}{0.8pt} % Rahmenstärke

    \fbox{%
      \begin{minipage}[t]{\linewidth}
        \renewcommand{\arraystretch}{1.25}
        \begin{tabular}{@{}p{0.24\linewidth}@{\hspace{10pt}}|@{\hspace{10pt}}p{0.70\linewidth}@{}}

          % --- LEFT: LOGO BOX ---
          \begin{minipage}[c][3.4cm][c]{\linewidth}
            \centering
            \includegraphics[width=0.92\linewidth]{C:/Users/TKast/OneDrive/Dokumente/GitHub/SchulaufgabenScript/logos/krs.png}
          \end{minipage}
          &

          % --- RIGHT: HEADER TABLE ---
          \begin{minipage}[t]{\linewidth}
            \renewcommand{\arraystretch}{1.25}
            \begin{tabular}{|p{0.44\linewidth}|p{0.26\linewidth}|p{0.26\linewidth}|}
              \hline
              \textbf{1. Schulaufgabe} & Fach: Mathematik & Datum: 2024-11-26 \\ \hline
          Name, Vorname: \linefield{7cm} & Klasse: \linefield{3cm} & Nr./Note: \linefield{3cm} \\ \hline
          \multicolumn{3}{|l|}{Zeitzuschlag: \linefield{12cm}} \\ \hline
          \multicolumn{3}{|l|}{Unterschrift Erziehungsberechtigte: \linefield{12cm}} \\ \hline
          \multicolumn{3}{|l|}{LRSt \checkbox \quad IRSt \checkbox \quad ILSt \checkbox} \\ \hline
            \end{tabular}
          \end{minipage}

        \end{tabular}
      \end{minipage}
    }

    \vspace{10pt}
    
\tasktitle{Aufgabe 1: Fehlerbilder zu linearen Funktionen}{4.5}
Verbinde jede Zeichnung mit der dazu passenden Beschreibung des Fehlers. Fehlt eine passende Beschreibung, formuliere sie im letzten Kasten.

\begin{center}\includegraphics[width=\linewidth]{C:/Users/TKast/OneDrive/Dokumente/GitHub/SchulaufgabenScript/tasks/lineare_funktionen_fehlerbilder_01/images/layout.png}\end{center}


\tasktitle{Aufgabe 2: Gerade g₁: Punktprobe, fehlende Koordinate, Nullstelle, Parallelität}{8.0}
Gegeben ist die Gerade $g_1: y = -\frac{2}{3}x + 2$.

\begin{enumerate}[label=\alph*)]
\item Prüfe rechnerisch, ob der Punkt $P\,(-1{,}5\mid 3)$ auf der Geraden $g_1$ liegt.
\begin{center}
\begin{tikzpicture}
  \draw[step=5mm, very thin] (0,0) grid ({\linewidth}, {4.0cm});
  \draw[line width=0.4pt] (0,0) rectangle ({\linewidth}, {4.0cm});
\end{tikzpicture}
\end{center}

\item Der Punkt $B\,(x\mid \frac{1}{3})$ liegt auf der Geraden $g_1$. Berechne die fehlende Koordinate.
\begin{center}
\begin{tikzpicture}
  \draw[step=5mm, very thin] (0,0) grid ({\linewidth}, {4.0cm});
  \draw[line width=0.4pt] (0,0) rectangle ({\linewidth}, {4.0cm});
\end{tikzpicture}
\end{center}

\item Berechne die Nullstelle der Geraden $g_1$.
\begin{center}
\begin{tikzpicture}
  \draw[step=5mm, very thin] (0,0) grid ({\linewidth}, {4.0cm});
  \draw[line width=0.4pt] (0,0) rectangle ({\linewidth}, {4.0cm});
\end{tikzpicture}
\end{center}

\item Prüfe rechnerisch, ob die Gerade $g_2$ mit der Gleichung $x + \frac{1}{3}y - 2 = 0$ parallel zur Geraden $g_1$ verläuft.
\begin{center}
\begin{tikzpicture}
  \draw[step=5mm, very thin] (0,0) grid ({\linewidth}, {5.0cm});
  \draw[line width=0.4pt] (0,0) rectangle ({\linewidth}, {5.0cm});
\end{tikzpicture}
\end{center}

\end{enumerate}
\vspace{6pt}\hrule\vspace{6pt}
\par\textbf{Gesamtpunkte:} 12.5\par
\end{document}